% SEGMENT OLP-0104-S01
% Part: first-order-logic
% Chapter: tableaux
% Section: proving-things-quant

\documentclass[../../../include/open-logic-section]{subfiles}

\begin{document}

\olfileid{fol}{tab}{prq}

\olsection{அளவையடைகளுடன் \usetoken{P}{tableau}}

% SEGMENT OLP-0104-S02
\begin{ex}
அளவையடைகளைக் கையாளும்போது தனிமாறி நிபந்தனையை மீறாமல் இருப்பதை உறுதிசெய்ய
வேண்டும்; இதற்காகச் சில அனுமானங்களைச் செயல்படுத்தும் வரிசையை மாற்றித் தேர்ந்து
பார்க்க வேண்டியிருக்கும். பொதுவாகத் தனிமாறி நிபந்தனைக்குட்பட்ட விதிகளை முதலில்
செய்துமுடிக்க முயல்வது உதவும் (நிறைவுற்ற !!{tableau}வில் அவை மேலே அமையும்).

!!{sentence} $\lexists[x][\lnot !A(x)] \lif \lnot \lforall[x][!A(x)]$
க்கான !!a{tableau}வை எவ்வாறு கொடுப்போம் என்று பார்ப்போம். வழக்கம்போல எடுகோளைப்
பதிவு செய்தே தொடங்குகிறோம்:
\begin{oltableau}
[\sFmla{\False}{\lexists[x][\lnot \formula{A}(x)]\lif \lnot
    \lforall[x][\formula{A}(x)]}, just=\TAss]
\end{oltableau}
!!{main operator}~$\lif$ என்பதால் $\TRule{\False}{\lif}$ ஐப் பயன்படுத்துகிறோம்:
\begin{oltableau}
[\sFmla{\False}{\lexists[x][\lnot \formula{A}(x)]\lif \lnot
    \lforall[x][\formula{A}(x)]}, just=\TAss, checked
  [\sFmla{\True}{\lexists[x][\lnot \formula{A}(x)]},
    just={\TRule{\False}{\lif}[1]}
    [\sFmla{\False}{\lnot\lforall[x][\formula{A}(x)]},
      just={\TRule{\False}{\lif}[1]}]
  ]
]
\end{oltableau}

% SEGMENT OLP-0104-S03
அடுத்து கையாள வேண்டியது வரி~$2$. $\TRule{\True}{\lexists}$ ஐப்
பயன்படுத்துகிறோம். இதற்கு ஒரு புதிய !!{constant} தேவை; இதுவரை எந்த
!!{constant}s உம் இடம்பெறாததால், எதை வேண்டுமானாலும் தேர்ந்தெடுக்கலாம்;
எடுத்துக்காட்டாக~$a$.
\begin{oltableau}
[\sFmla{\False}{\lexists[x][\lnot \formula{A}(x)]\lif \lnot
    \lforall[x][\formula{A}(x)]}, just=\TAss, checked
  [\sFmla{\True}{\lexists[x][\lnot \formula{A}(x)]},
    just={\TRule{\False}{\lif}[1]}, checked
    [\sFmla{\False}{\lnot\lforall[x][\formula{A}(x)]},
      just={\TRule{\False}{\lif}[1]}
      [\sFmla{\True}{\lnot\formula{A}(a)}, just={\TRule{\True}{\lexists}[2]}]
    ]
  ]
]
\end{oltableau}

% SEGMENT OLP-0104-S04
இப்போது வரி~$3$ க்கு $\TRule{\False}{\lnot}$ ஐப் பயன்படுத்துகிறோம்:
\begin{oltableau}
[\sFmla{\False}{\lexists[x][\lnot \formula{A}(x)]\lif \lnot
    \lforall[x][\formula{A}(x)]}, just=\TAss, checked
  [\sFmla{\True}{\lexists[x][\lnot \formula{A}(x)]},
    just={\TRule{\False}{\lif}[1]}, checked
    [\sFmla{\False}{\lnot\lforall[x][\formula{A}(x)]},
      just={\TRule{\False}{\lif}[1]}, checked
      [\sFmla{\True}{\lnot\formula{A}(a)}, just={\TRule{\True}{\lexists}[2]}
        [\sFmla{\True}{\lforall[x][\formula{A}(x)]},
          just = {\TRule{\False}{\lnot}[3]}
        ]
      ]
    ]
  ]
]
\end{oltableau}

% SEGMENT OLP-0104-S05
வரி~$4$ க்கு $\TRule{\True}{\lnot}$ ஐயும், அதைத் தொடர்ந்து வரி~$5$ க்கு
$\TRule{\True}{\lforall}$ ஐயும் பயன்படுத்தி மூடிய !!a{tableau}வைப் பெறுகிறோம்.
\begin{oltableau}
[\sFmla{\False}{\lexists[x][\lnot \formula{A}(x)]\lif \lnot
    \lforall[x][\formula{A}(x)]}, just=\TAss, checked
  [\sFmla{\True}{\lexists[x][\lnot \formula{A}(x)]},
    just={\TRule{\False}{\lif}[1]}, checked
    [\sFmla{\False}{\lnot\lforall[x][\formula{A}(x)]},
      just={\TRule{\False}{\lif}[1]}, checked
      [\sFmla{\True}{\lnot\formula{A}(a)}, just={\TRule{\True}{\lexists}[2]}
        [\sFmla{\True}{\lforall[x][\formula{A}(x)]},
          just = {\TRule{\False}{\lnot}[3]}
          [\sFmla{\False}{\formula{A}(a)}, just={\TRule{\True}{\lnot}[4]}
            [\sFmla{\True}{\formula{A}(a)},
              just={\TRule{\True}{\lforall}[5]}, close
            ]
          ]
        ]
      ]
    ]
  ]
]
\end{oltableau}
\end{ex}

% SEGMENT OLP-0104-S06
\begin{ex}
பின்வரும் கணத்திற்கான !!a{tableau}வை எவ்வாறு கொடுப்போம் என்று பார்ப்போம்:
\[
\sFmla{\False}{\lexists[x][!C(x,b)]},
\sFmla{\True}{\lexists[x][(!A(x) \land !B(x))]},
\sFmla{\True}{\lforall[x][(!B(x) \lif !C(x,b))]}.
\]
வழக்கம்போல எடுகோள்களுடன் தொடங்குகிறோம்:
\begin{oltableau}
  [\sFmla{\False}{\lexists[x][\formula{C}(x,b)]}, just=\TAss
    [\sFmla{\True}{\lexists[x][(\formula{A}(x) \land \formula{B}(x))]},
      just=\TAss
      [\sFmla{\True}{\lforall[x][(\formula{B}(x) \lif \formula{C}(x,b))]},
        just=\TAss]
    ]
  ]
\end{oltableau}

% SEGMENT OLP-0104-S07
தனிமாறி நிபந்தனையுள்ள விதியை எப்போதும் முதலில் பயன்படுத்த வேண்டும்; இங்கு
அது வரி~$2$ க்கான $\TRule{\True}{\lexists}$. எடுகோள்களில்
!!{constant}~$b$ இருப்பதால் வேறொன்றைப் பயன்படுத்த வேண்டும்; மீண்டும்~$a$ ஐத்
தேர்ந்தெடுப்போம்.
\begin{oltableau}
  [\sFmla{\False}{\lexists[x][\formula{C}(x,b)]}, just=\TAss
    [\sFmla{\True}{\lexists[x][(\formula{A}(x) \land \formula{B}(x))]},
      just=\TAss, checked
      [\sFmla{\True}{\lforall[x][(\formula{B}(x) \lif \formula{C}(x,b))]},
        just=\TAss
        [\sFmla{\True}{\formula{A}(a) \land \formula{B}(a)},
          just={\TRule{\True}{\lexists}[2]}]
      ]
    ]
  ]
\end{oltableau}

% SEGMENT OLP-0104-S08
இப்போது வரி~$1$ க்கு $\TRule{\False}{\lexists}$ ஐயோ, வரி~$3$ க்கு
$\TRule{\True}{\lforall}$ ஐயோ பயன்படுத்தினால், $x$ க்குப் பதிலாக எந்தச்
சொல்~$t$ ஐ இடுவது என்று தீர்மானிக்க வேண்டும். இவ்விதிகளுக்குத் தனிமாறி
நிபந்தனை இல்லாததால் விரும்பிய சொல்லைத் தேர்ந்தெடுக்கலாம். சில நேர்வுகளில்
வெவ்வேறு~$t$ களுடன் விதியைப் பலமுறை பயன்படுத்தவும் வேண்டியிருக்கலாம். ஆனால்
பொதுவிதியாக, !!{tableau}வில் ஏற்கெனவே இடம்பெறும் சொற்களில் ஒன்றை---இங்கு
$a$, $b$---தேர்ந்தெடுப்பது பயனளிக்கும்; இந்நேர்வில் $a$ ஒரு மூடிய கிளையைத்
தர அதிக வாய்ப்புள்ளது என்று ஊகிக்கலாம்.
\begin{oltableau}
  [\sFmla{\False}{\lexists[x][\formula{C}(x,b)]}, just=\TAss
    [\sFmla{\True}{\lexists[x][(\formula{A}(x) \land \formula{B}(x))]},
      just=\TAss, checked
      [\sFmla{\True}{\lforall[x][(\formula{B}(x) \lif \formula{C}(x,b))]}, just=\TAss
        [\sFmla{\True}{\formula{A}(a) \land \formula{B}(a)}, just={\TRule{\True}{\lexists}[2]}
          [\sFmla{\False}{\formula{C}(a,b)}, just={\TRule{\False}{\lexists}[1]}
            [\sFmla{\True}{\formula{B}(a) \lif \formula{C}(a,b)},
              just={\TRule{\True}{\lforall}[3]}
            ]
          ]
        ]
      ]
    ]
  ]
\end{oltableau}

% SEGMENT OLP-0104-S09
வரி $1$, $3$ இலுள்ள !!{signed formula}s ஐச் சரிக்குறியிடுவதில்லை; ஏனெனில்
அவற்றை மீண்டும் பயன்படுத்த வேண்டியிருக்கலாம். இப்போது வரி~$4$ க்கு
$\TRule{\True}{\land}$ ஐப் பயன்படுத்துக:
\begin{oltableau}
  [\sFmla{\False}{\lexists[x][\formula{C}(x,b)]}, just=\TAss
    [\sFmla{\True}{\lexists[x][(\formula{A}(x) \land \formula{B}(x))]},
      just=\TAss, checked
      [\sFmla{\True}{\lforall[x][(\formula{B}(x) \lif \formula{C}(x,b))]}, just=\TAss
        [\sFmla{\True}{\formula{A}(a) \land \formula{B}(a)},
          just={\TRule{\True}{\lexists}[2]}, checked
          [\sFmla{\False}{\formula{C}(a,b)}, just={\TRule{\False}{\lexists}[1]}
            [\sFmla{\True}{\formula{B}(a) \lif \formula{C}(a,b)},
              just={\TRule{\True}{\lforall}[3]}
              [\sFmla{\True}{\formula{A}(a)}, just={\TRule{\True}{\land}[4]}
                [\sFmla{\True}{\formula{B}(a)}, just={\TRule{\True}{\land}[4]}
                ]
              ]
            ]
          ]
        ]
      ]
    ]
  ]
\end{oltableau}

% SEGMENT OLP-0104-S10
இப்போது வரி~$6$ க்கு $\TRule{\True}{\lif}$ ஐப் பயன்படுத்தினால் !!{tableau}
மூடுகிறது:
\begin{oltableau}
  [\sFmla{\False}{\lexists[x][\formula{C}(x,b)]}, just=\TAss
    [\sFmla{\True}{\lexists[x][(\formula{A}(x) \land \formula{B}(x))]},
      just=\TAss, checked
      [\sFmla{\True}{\lforall[x][(\formula{B}(x) \lif \formula{C}(x,b))]},
        just=\TAss
        [\sFmla{\True}{\formula{A}(a) \land \formula{B}(a)},
          just={\TRule{\True}{\lexists}[2]}, checked
          [\sFmla{\False}{\formula{C}(a,b)},
            just={\TRule{\False}{\lexists}[1]}
            [\sFmla{\True}{\formula{B}(a) \lif \formula{C}(a,b)},
              just={\TRule{\True}{\lforall}[3]}, checked
              [\sFmla{\True}{\formula{A}(a)},
                just={\TRule{\True}{\land}[4]}
                [\sFmla{\True}{\formula{B}(a)},
                  just={\TRule{\True}{\land}[4]}
                  [\sFmla{\False}{\formula{B}(a)},
                    just={\TRule{\True}{\lif}[6]},close
                  ]
                  [\sFmla{\True}{\formula{C}(a,b)},
                    just={\TRule{\True}{\lif}[6]},close
                  ]
                ]
              ]
            ]
          ]
        ]
      ]
    ]
  ]
\end{oltableau}


\end{ex}

% SEGMENT OLP-0104-S11
\begin{ex}
பின்வரும் கணத்திற்கான !!a{tableau}வை அமைக்கிறோம்:
\[
\sFmla{\True}{\lforall[x][!A(x)]}, \sFmla{\True}{\lforall[x][!A(x)] 
  \lif \lexists[y][!B(y)]}, \sFmla{\True}{\lnot\lexists[y][!B(y)]}.
\]
வழக்கம்போல எடுகோள்களை எழுதுகிறோம்:
\begin{oltableau}
  [\sFmla{\True}{\lforall[x][\formula{A}(x)]}, just=\TAss
    [\sFmla{\True}{\lforall[x][\formula{A}(x)] \lif
        \lexists[y][\formula{B}(y)]}, just=\TAss
      [\sFmla{\True}{\lnot\lexists[y][\formula{B}(y)]}, just=\TAss
      ]
    ]
  ]
\end{oltableau}

% SEGMENT OLP-0104-S12
வரி~$3$ க்கு $\TRule{\True}{\lnot}$ விதியைப் பயன்படுத்தித் தொடங்குகிறோம்.
``தனிமாறி நிபந்தனையுள்ள விதிகளை எப்போதும் முதலில் பயன்படுத்துக'' என்ற விதியின்
ஒரு பின்விளைவு, ``தனிமாறி நிபந்தனையில்லாத அளவையடை விதிகள் தேவைப்படும் வரை
அவற்றைப் பயன்படுத்துவதை ஒத்திவைக்கவும்'' என்பதாகும். கிளைப்பை விளைவிக்கும்
விதிகளையும் ஒத்திவைக்கவும்.
\begin{oltableau}
  [\sFmla{\True}{\lforall[x][\formula{A}(x)]}, just=\TAss
    [\sFmla{\True}{\lforall[x][\formula{A}(x)] \lif
        \lexists[y][\formula{B}(y)]}, just=\TAss
      [\sFmla{\True}{\lnot\lexists[y][\formula{B}(y)]}, just=\TAss, checked
        [\sFmla{\False}{\lexists[y][\formula{B}(y)]}, just={\TRule{\True}{\lnot}[3]}]
      ]
    ]
  ]
\end{oltableau}

% SEGMENT OLP-0104-S13
புதிய வரி~$4$ க்குத் தனிமாறி நிபந்தனையில்லாத அளவையடை விதியான
$\TRule{\False}{\lexists}$ தேவை. ஆகவே இதை ஒத்திவைத்து, வரி~$2$ இல்
$\TRule{\True}{\lif}$ ஐப் பயன்படுத்துகிறோம்.
\begin{oltableau}
  [\sFmla{\True}{\lforall[x][\formula{A}(x)]}, just=\TAss
    [\sFmla{\True}{\lforall[x][\formula{A}(x)] \lif
        \lexists[y][\formula{B}(y)]}, just=\TAss, checked
      [\sFmla{\True}{\lnot\lexists[y][\formula{B}(y)]}, just=\TAss, checked
        [\sFmla{\False}{\lexists[y][\formula{B}(y)]},
          just={\TRule{\True}{\lnot}[3]},
          [\sFmla{\False}{\lforall[x][\formula{A}(x)]}, just={\TRule{\True}{\lif}[2]}]
          [\sFmla{\True}{\lexists[y][\formula{B}(y)]}, just={\TRule{\True}{\lif}[2]}]
        ]
      ]
    ]
  ]
\end{oltableau}

% SEGMENT OLP-0104-S14
புதிய !!{signed formula}s இரண்டிற்கும் தனிமாறி நிபந்தனையுள்ள விதிகள் தேவை;
ஆகவே அவற்றை அடுத்து பயன்படுத்த வேண்டும்:
\begin{oltableau}
  [\sFmla{\True}{\lforall[x][\formula{A}(x)]}, just=\TAss
    [\sFmla{\True}{\lforall[x][\formula{A}(x)] \lif
        \lexists[y][\formula{B}(y)]}, just=\TAss, checked
      [\sFmla{\True}{\lnot\lexists[y][\formula{B}(y)]}, just=\TAss, checked
        [\sFmla{\False}{\lexists[y][\formula{B}(y)]},
          just={\TRule{\True}{\lnot}[3]}
          [\sFmla{\False}{\lforall[x][\formula{A}(x)]}, just={\TRule{\True}{\lif}[2]},checked
            [\sFmla{\False}{\formula{A}(b)}, just={\TRule{\False}{\lforall}[5]}]
          ]
          [\sFmla{\True}{\lexists[y][\formula{B}(y)]}, just={\TRule{\True}{\lif}[2]},checked
            [\sFmla{\True}{\formula{B}(c)}, just={\TRule{\True}{\lexists}[5]}]
          ]
        ]
      ]
    ]
  ]
\end{oltableau}

% SEGMENT OLP-0104-S15
கிளைகளை மூட, வரி $1$, $3$ இலுள்ள !!{signed formula}s ஐப் பயன்படுத்த வேண்டும்.
அவற்றிற்குரிய விதிகளுக்கு (\TRule{\True}{\lforall},
\TRule{\False}{\lexists}) தனிமாறி நிபந்தனைகள் இல்லை; எனவே பொருத்தமான எந்தச்
சொற்களையும் தேர்ந்தெடுக்கலாம். இங்கு அவை முறையே $b$, $c$.
\begin{oltableau}
  [\sFmla{\True}{\lforall[x][\formula{A}(x)]}, just=\TAss
    [\sFmla{\True}{\lforall[x][\formula{A}(x)] \lif
        \lexists[y][\formula{B}(y)]}, just=\TAss, checked
      [\sFmla{\True}{\lnot\lexists[y][\formula{B}(y)]}, just=\TAss, checked
        [\sFmla{\False}{\lexists[y][\formula{B}(y)]},
          just={\TRule{\True}{\lnot}[3]}
          [\sFmla{\False}{\lforall[x][\formula{A}(x)]},
            just={\TRule{\True}{\lif}[2]},checked
            [\sFmla{\False}{\formula{A}(b)},
              just={\TRule{\False}{\lforall}[5]}
              [\sFmla{\True}{\formula{A}(b)},
                just={\TRule{\True}{\lforall}[1]},close
              ]
            ]
          ]
          [\sFmla{\True}{\lexists[y][\formula{B}(y)]},
            just={\TRule{\True}{\lif}[2]},checked
            [\sFmla{\True}{\formula{B}(c)},
              just={\TRule{\True}{\lexists}[5]}
              [\sFmla{\False}{\formula{B}(c)},
                just={\TRule{\False}{\lexists}[4]},close
              ]
            ]
          ]
        ]
      ]
    ]
  ]
\end{oltableau}
\end{ex}

% SEGMENT OLP-0104-S16
\begin{prob}
பின்வருவனவற்றிற்கான மூடிய !!{tableau}s ஐக் கொடுக்கவும்:
\begin{enumerate}
\item $\sFmla{\False}{(\lforall[x][!A(x)] \land \lforall[y][!B(y)])
\lif \lforall[z][(!A(z) \land !B(z))]}$.
\item $\sFmla{\False}{(\lexists[x][!A(x)] \lor \lexists[y][!B(y)])
\lif \lexists[z][(!A(z) \lor !B(z))]}$.
\item $\sFmla{\True}{\lforall[x][(!A(x) \lif !B)]},
\sFmla{\False}{\lexists[y][!A(y)] \lif !B}$.
\item $\sFmla{\True}{\lforall[x][\lnot !A(x)]},
\sFmla{\False}{\lnot\lexists[x][!A(x)]}$.
\item $\sFmla{\False}{\lnot\lexists[x][!A(x)] \lif \lforall[x][\lnot
!A(x)]}$.
\item $\sFmla{\False}{\lnot\lexists[x][\lforall[y][((!A(x,y) \lif
\lnot !A(y,y)) \land (\lnot !A(y,y) \lif !A(x,y)))]]}$.
\end{enumerate}
\end{prob}

% SEGMENT OLP-0104-S17
\begin{prob}
பின்வருவனவற்றிற்கான மூடிய !!{tableau}s ஐக் கொடுக்கவும்:
\begin{enumerate}
\item $\sFmla{\False}{\lnot\lforall[x][!A(x)] \lif \lexists[x][\lnot!A(x)]}$.
\item $\sFmla{\True}{(\lforall[x][!A(x)] \lif !B)}, \sFmla{\False}{\lexists[y][(!A(y) \lif !B)]}$.
\item $\sFmla{\False}{\lexists[x][(!A(x) \lif \lforall[y][!A(y)])]}$.
\end{enumerate}
\end{prob}

\end{document}
